\chapter{Hardware wiring}
\label{ch:wiring}
\index{hardware}\index{wiring}

This chapter is an outline. It records what the software requires of the hardware,
taken from the code and from the names in the TDT circuit files. The physical
connections have not been checked on the rig, and need to be written up by someone who
has access to it.

\section{Equipment}
\index{equipment}

\begin{center}
\begin{tabularx}{\textwidth}{@{}lX@{}}
\toprule
Device & Role \\
\midrule
NI 6371\index{NI 6371} & Stimulus output at 500~kHz, $\pm$10~V, 16 bits \\
TDT RP2.1\index{RP2.1} & Records the ABR and the microphone; generates the trigger \\
TDT PA5\index{PA5} & Programmable attenuator, controlled over USB \\
Speaker driver & Drives the speaker. The setting is confirmed by hand before
  calibration \\
Speaker & MF1, EC1 or ES1 family (Table~\ref{tab:speakers}) \\
Microphone and preamplifier & Used for calibration (Table~\ref{tab:microphones}) \\
Sound level calibrator & 94~dB at 1~kHz \\
Biological amplifier & Amplifies the electrode signal, with a gain of 10,000 \\
TDT RZ5D\index{RZ5D} & Not used in this configuration. It could be used for acquisition
  if some software problems were resolved, but it cannot produce the output frequency
  needed for tones above 90~kHz \\
\bottomrule
\end{tabularx}
\end{center}

\needsinput{Exact models of the NI card, speaker driver, microphone preamplifier,
calibrator and biological amplifier, with the amplifier's filter settings.}

\section{Signal paths}
\index{signal paths}

\begin{figure}[htbp]
\centering
\begin{tikzpicture}[
    node distance=6mm and 10mm,
    dev/.style={draw=uncblue!80!black,fill=uncblue!12,rounded corners=2pt,thick,
                minimum width=26mm,minimum height=10mm,align=center,font=\small},
    ext/.style={dev,fill=gray!12,draw=gray!70},
    sig/.style={-{Latex[length=2.5mm]},thick},
    trg/.style={sig,dashed,draw=orange!70!black},
    pin/.style={font=\scriptsize,midway,above}]
  \node[dev] (rp) {TDT RP2.1};
  \node[dev,right=30mm of rp] (ni) {NI card};
  \node[dev,right=18mm of ni] (pa) {PA5};
  \node[ext,below=12mm of pa] (drv) {Speaker\\driver};
  \node[ext,left=18mm of drv] (spk) {Speaker};
  \node[ext,left=18mm of spk] (animal) {Animal or\\microphone};
  \node[ext,below=of animal] (amp) {Amplifier /\\mic preamp};
  \draw[trg] (rp) -- node[pin]{trigger to PFI0} (ni);
  \draw[sig] (ni) -- node[pin]{AO0} (pa);
  \draw[sig] (pa) -- (drv);
  \draw[sig] (drv) -- (spk);
  \draw[sig] (spk) -- node[pin]{sound} (animal);
  \draw[sig] (animal) -- (amp);
  \draw[sig] (amp.west) -- ++(-20mm,0) |- node[pos=0.25,left,font=\scriptsize,align=right]{In1 (ABR)\\In2 (mic)} ($(rp.south west)!0.25!(rp.north west)$);
\end{tikzpicture}
\caption{Signal paths inferred from the software. The dashed line is the trigger.}
\label{fig:signal-paths}
\end{figure}

\begin{table}[htbp]
\centering
\caption{Connections and settings that the software depends on.}
\label{tab:connections}
\begin{tabularx}{\textwidth}{@{}lXl@{}}
\toprule
Item & Requirement & Status \\
\midrule
Stimulus output & NI analog output \texttt{ao0}, $\pm$10~V & From code \\
Trigger & NI \texttt{PFI0}, rising edge, starts the stimulus. For ABR recording the
  device name in this path is fixed as \texttt{Dev1} & From code \\
Trigger source & A digital output of the RP2.1 (a BitOut block labeled ``NI Trigger Pulse
  Generator'' in both circuits) & Port to verify \\
ABR input & RP2.1 In1, stored in buffer \texttt{sampled\_wave}, at
  24414.0625~Hz & Confirmed \\
Microphone input & RP2.1 In2, read from \texttt{Data\_out2}, at
  195312.5~Hz & Confirmed \\
Attenuator & PA5 number 1, over USB. The rig has two PA5 units, but the code controls
  only this one & From code \\
RP2.1 connection & USB interface & From code \\
Amplifier gain & 10,000. The software cannot read the gain from the amplifier, so the
  data viewer uses this entered value & Confirmed \\
Mic preamp gain & 20~dB by default; must match the microphone calibration & Verify \\
Speaker driver & $-$6~dB by default & Verify \\
\bottomrule
\end{tabularx}
\end{table}

\figorplaceholder{tdt_circuit_abr}{The RP2.1 circuit used for ABR recording,
\file{tdt/abrs\_v2.rcx}.}{Export from RPvdsEx of tdt/abrs\_v2.rcx.}

\figorplaceholder{tdt_circuit_mic}{The RP2.1 circuit used for calibration,
\file{tdt/mic\_record.rcx}.}{Export from RPvdsEx of tdt/mic\_record.rcx.}

\section{Connections}

\needsinput{A cable-by-cable list: from, to, cable type and any labels. Include the
RP2.1 output that carries the trigger, which is not certain from the software.}

\figorplaceholder{hw_rack_front}{The equipment rack, front.}{Photograph of the front of
the rack, with each device labeled.}

\figorplaceholder{hw_rack_back}{The equipment rack, rear connections.}{Photograph of the
rear connections, with the trigger, output and input cables identified.}

\section{Sound delivery and recording setup}

\needsinput{Speaker position relative to the animal, free-field or closed-field use,
electrode positions and types, grounding, and steps taken to reduce electrical noise
(for example, the heating pad). Include the switch settings to be used during
recording.}

\figorplaceholder{hw_speaker_setup}{Speaker and animal position.}{Photograph of the
speaker and the animal (or a stand-in), showing distance and angle.}

\figorplaceholder{hw_electrodes}{Electrode placement.}{Photograph or drawing of the
electrode positions.}

\figorplaceholder{hw_calibrator}{The microphone in the calibrator.}{Photograph of the
microphone seated in the sound level calibrator.}

\section{Power-up and checks}
\index{power-up}

\needsinput{The order in which equipment is switched on, and a short check that the
sound path works before starting. The output stage can make noise if the attenuators
are not at maximum.}

\figorplaceholder[0.7\textwidth]{ni_max_device}{The NI device name in the NI-DAQmx
utility, as used for \key{NI\_device}.}{Screen capture of NI MAX, showing the device
name (Dev1).}
